\section{An exponential lower bound at the mean-field exponent}
\label{sec:critical-depth-lower}

At the critical exponent $\gamma=1$, a different argument gives an
exponential lower bound for the same teacher family and the same fixed
nonatomic input distribution $D$. The width tends to infinity with
teacher depth and the realized sample fixed. Only finite-dimensional
large deviations are needed: apply them separately to the training
sample and to one test pair, then integrate over the test pairing.

\begin{theorem}[Critical regression lower bound]
\label{thm:critical-regression-depth-lower}
Use the distribution $D$ and targets $f_h(x)=F^h(x_1)$ from
Theorem~\ref{thm:strong}. There are universal constants $c>0$ and $h_0$
such that the following holds for every $h\ge h_0$ and integer
$1\le n\le e^{ch}$. Draw $n$ independent inputs from $D$ and observe
their exact teacher outputs. Use the original Gaussian prior at depth
$L=h+1$, with first-layer entry variance $1/2$, internal entry variance
$1/N$, and readout entry variance $1/N^2$. Condition solely on
\[
 |z_w(X_i)-f_h(X_i)|<\rho,\qquad \rho=1/8.
\]
With sample probability at least $49/50$, the posterior satisfies
\[
 \liminf_{N\to\infty}
 \mathbb E_{w\mid S}D\{|z_w-f_h|>1/4\}\ge\frac12.
\]
Moreover, writing
$p_N(x)=\Pr_{w\mid S}\{|z_w(x)-f_h(x)|\le1/4\}$, the same samples obey
\[
 \liminf_{N\to\infty}D\{x:p_N(x)<2/3\}\ge\frac12.
\]
The teacher still has width three and weights bounded by thirty.
There is no global output-margin condition. These statements concern
the actual posterior, including mixtures, and give no width threshold
uniform in growing $h$.
\end{theorem}

\subsection{Finite-domain rates and a metric bound}

The rate below measures the exponential prior cost of approximating
prescribed outputs on finitely many inputs. Its hidden-layer terms are
relative entropies, and its final term is the Gaussian readout cost.
We first show that a small total cost limits the output difference
between nearby inputs.

We include the large-deviation argument needed below. It works for any
fixed finite input set $X=\{x^1,\ldots,x^m\}\subset S^1$. For a feature
law $\nu$ on $E=[-1,1]^m$, write
$Q(\nu)=\int hh^T\,d\nu(h)$. Let $\mathsf M_Q$ be the law of
coordinatewise tanh applied to $\mathcal N(0,Q)$, and put
\[
 Q_0(x,y)=\frac{1+x\cdot y}{2},\qquad
 \mathcal G_Q(z)=
 \begin{cases}\frac12z^TQ^+z,&z\in\operatorname{range}Q,\\
 +\infty,&\text{otherwise}.
 \end{cases}
\]
Here and below singular covariance matrices are allowed.

\begin{lemma}[Finite-domain critical large deviations]
\label{lem:critical-finite-ldp-depth}
The output vector $Z_N=(z_w(x))_{x\in X}$ obeys a large-deviation
principle with speed $N$ and good rate
\[
 I_L(z)=\inf_{\nu_1,\ldots,\nu_{L-1}}
 \left\{\sum_{s=1}^{L-1}D(\nu_s\Vert\mathsf M_{Q_{s-1}})
                   +\mathcal G_{Q_{L-1}}(z)\right\},
 \qquad Q_s=Q(\nu_s).
\]
For every finite-cost path of total cost $A$ in this formula,
\begin{equation}\label{eq:critical-metric-bound-depth}
 |z(x)-z(y)|\le\sqrt A\,|x-y|
                         \exp\!\bigl(\sqrt{2(L-1)A}\bigr)
 \qquad(x,y\in X).
\end{equation}
\end{lemma}

\begin{proof}
Let $\Gamma_m=\mathcal N(0,I_m)$, and take independent empirical
measures $\eta_{s,N}$ of $N$ samples from $\Gamma_m$, for
$1\le s<L$. Independently take $V_N=G/\sqrt N$, $G\sim\Gamma_m$.
For $\mathsf T_Q(u)=\tanh(Q^{1/2}u)$, recursively define
\[
 \nu_{1,N}=(\mathsf T_{Q_0})_\#\eta_{1,N},\qquad
 \nu_{s,N}=(\mathsf T_{Q(\nu_{s-1,N})})_\#\eta_{s,N}.
\]
These are exactly the network's empirical hidden-feature laws on $X$:
conditionally on the preceding layer, its next rows are independent
Gaussian fields with that empirical second moment. The final output
has exactly the law $Q(\nu_{L-1,N})^{1/2}V_N$.

Sanov's theorem and the Gaussian large-deviation principle give joint
good rate $\sum_s D(\eta_s\Vert\Gamma_m)+\|v\|^2/2$.
The recursive map just described is continuous in the weak topology
for the measures and the ordinary topology for $v$. Indeed, matrix
square roots are continuous; the bounded continuous test functions
composed with $\mathsf T_Q$ converge uniformly on compact sets when
$Q$ varies; and weakly convergent measures form a tight sequence.
Second moments on the compact feature space $E$ are also continuous.
The contraction principle therefore applies. The entropy-contraction
identity
\[
 \inf_{\eta:T_\#\eta=\nu}D(\eta\Vert\Gamma_m)
       =D(\nu\Vert T_\#\Gamma_m)
\]
follows by data processing and, for the reverse inequality, the lift
$d\eta=q(T(u))\,d\Gamma_m(u)$ when
$q=d\nu/d(T_\#\Gamma_m)$. Finally
$\inf_{Q^{1/2}v=z}\|v\|^2/2=\mathcal G_Q(z)$.
These identities give the displayed rate, including the singular cases.
This is an application of the standard Sanov and contraction theorems
\cite{dembozeitouni}.

For the metric bound, write $D_s$ for each hidden-layer entropy and
$r$ for the readout cost. Inverse tanh lifts a finite-entropy feature
law to a law $\eta$ with the same entropy relative to its reference
Gaussian. The Gaussian quadratic transport inequality
\cite{talagrand1996}, applied after whitening on the covariance range,
gives a coupling $(U,G)$ satisfying
\[
 \mathbb E\|Q^{+1/2}(U-G)\|^2\le2D_s.
\]
For $b=e_x-e_y$, Cauchy--Schwarz and Minkowski imply
\[
 \|b^TU\|_{L^2}\le(1+\sqrt{2D_s})\sqrt{b^TQb}.
\]
Tanh is $1$-Lipschitz, so this is also an upper bound on the next
feature distance. The initial distance is $|x-y|/\sqrt2$.
The readout inequality is
$|b^Tz|\le\sqrt{2r}\sqrt{b^TQ_{L-1}b}$.
Multiply these inequalities and use
\[
 \sum_{s=1}^{L-1}\log(1+\sqrt{2D_s})
 \le\sqrt{2(L-1)\sum_sD_s}\le\sqrt{2(L-1)A},\qquad r\le A.
\]
This proves~\eqref{eq:critical-metric-bound-depth}.
\end{proof}

Each test pair $(X_+,X_-)$ from Lemma~\ref{lem:pairing} has
input distance at most $3\cdot3^{-h}$ and target gap at least one.
Predicting both to tolerance $1/4$ requires an output gap of at least
$1/2$. At $L=h+1$, the metric bound shows that every such output on
the two-point domain has
\begin{equation}\label{eq:critical-correct-pair-rate}
 I_L(z)\ge h/100
\end{equation}
for all sufficiently large $h$: a path of cost $A\le h/100$ would
instead give gap at most
$3\cdot3^{-h}\sqrt{h/100}\exp(h/\sqrt{50})$, which tends to zero.
The same conclusion holds after taking the infimum over paths.

\subsection{A low-rate training interpolator}

The training constraints, in contrast, can be met at a much smaller
rate. The next construction spends entropy only in the first layer and
in logarithmically many final layers. Most hidden layers follow the
uncontrolled Gaussian recursion, which reduces their variances while
retaining a weak signal correlated with the labels. The final controlled
layers amplify that signal.

\begin{lemma}[A first-feature mean seed]
\label{lem:critical-mean-seed-depth}
Suppose $t_1<\cdots<t_n$ belong to $[-2/3,2/3]$ and have
separation at least $\Delta\le1$. For any signs $y_i$ and any fixed
finite domain containing $x(t_i)$, a genuine first-feature law $\nu$
satisfies
\[
 D(\nu\Vert\mathsf M_{Q_0})\le D,\qquad
 \mathbb E_\nu H(x(t_i))=m_n y_i,\qquad m_n=\frac1{2n},
\]
where $1\le D\le C\Delta^{-2}[1+\log(n+1)]^2$ for a universal $C$.
\end{lemma}

\begin{proof}
Use the ideal step vectors $s_0=1$ and
$s_j(i)=\operatorname{sign}(t_i-(t_j+t_{j+1})/2)$, $1\le j<n$.
The inverse of this basis maps the maximum norm to coefficient
$\ell^1$ norm with norm at most $n$: the coefficients are successive
half-differences and the half-sum of the two endpoint values.

In first-layer weight space, shift the bias and first coordinate so
that a row's preactivation mean is $B(t-u_j)$, where $u_j$ is the
corresponding midpoint. For the constant vector use constant mean $B$.
The random preactivation variance is one on $S^1$. Since
\[
 1-\mathbb E\tanh(m+Z)\le2e^{-2m+2}\qquad(Z\sim\mathcal N(0,1)),
\]
choosing $B=C\Delta^{-1}[1+\log(n+1)]$ makes the resulting mean
vectors $v_j$ approximate $s_j$ to maximum error $1/(4n)$.
A Neumann series gives exact coefficients
$y=\sum_j c_jv_j$ with $\sum_j|c_j|\le4n/3$.
Mix the corresponding signed Gaussian shifts with weights
$m_n|c_j|$, and put the remaining probability on the unshifted
Gaussian law. Negating a shift negates its feature mean by oddness.
The desired mean identity follows. Each shift has Gaussian relative
entropy at most $CB^2$, and convexity and data processing give the
stated bound. These are laws of actual weight vectors, so their
feature means and covariances are defined on the entire finite domain,
not only on the sample.
\end{proof}

\begin{lemma}[Free cooling and late amplification]
\label{lem:critical-late-amplification-depth}
Suppose a first-feature law has entropy at most $D\ge1$ and mean
$a_0y_i$ at selected inputs, where $y_i\in\{-1,1\}$ and
$0<a_0\le1$. For every fixed $0<\eta<1$, there are constants
$C_1,C_2$, depending only on $\eta$, such that if
\[
 L\ge C_1[1+\log(D/a_0^2)+\log(L+1)],
\]
then the critical rate admits an output $z$ with
\[
 |z(x_i)-a y_i|<a\eta,
 \qquad I_L(z)\le C_2[1+\log(D/a_0^2)+\log(L+1)],
 \qquad a=3/5.
\]
\end{lemma}

\begin{proof}
Mix the supplied law with the reference first-feature law at weight
$1/D$. This costs at most one and yields a covariance $Q\succeq bb^T$
with $b(x_i)=a_0y_i/\sqrt D$. The vector $b$ is the full-domain mean
of the supplied law divided by $\sqrt D$.

Consider an uncontrolled transition with $Q\succeq bb^T$ and
$q_x=Q(x,x)\le1$. Write the reference Gaussian field as $bZ_0+G_R$.
Project its tanh features onto $Z_0$. Gaussian integration by parts
shows that the resulting covariance dominates $b'b'^T$, where
\[
 b'_x=b_xh(q_x),\qquad
 h(q)=\mathbb E\operatorname{sech}^2(\sqrt qZ)=1-F_0(q),\qquad
 F_0(q)=\mathbb E\tanh^2(\sqrt qZ).
\]
The new diagonal is $F_0(q_x)$. There are universal $c>0,C<\infty$
such that
\[
 F_0(q)\le q-cq^2,\qquad -\log h(q)\le Cq\quad(0\le q\le1).
\]
For example $c=1/100$ works. On $|u|\le1$,
$u^2-\tanh^2u\ge u^4/6$; integrate over $|Z|\le1$ to obtain the
first bound. For the second, $h(q)\ge h(1)\ge1/8$ and $F_0(q)\le q$.
It follows that after $H$ free layers the diagonal is at most
$(1+cH)^{-1}$ and the signed seed at every selected input is at least
\begin{equation}\label{eq:critical-free-seed}
 u_*:=e^{-8}(L+1)^{-800}a_0/\sqrt D
 \qquad(H\le L).
\end{equation}
Indeed the diagonal reciprocals grow by at least $c$ per step, and
summing $-\log h(q)\le8q$ proves this loose explicit bound.

Now suppose $Q\succeq bb^T$ and $y_ib_i\ge u>0$ on the sample.
Shift the reference Gaussian by $A b$. The entropy cost is at most
$A^2/2$, since $b^TQ^+b\le1$. The new mean $b'$ satisfies
\[
 y_i b'_i\ge c_0\min\{Au,1\},\qquad c_0=1/100,
\]
and the new covariance dominates $b'b'^T$. To check the inequality,
integrate the derivative
$\mathbb E\operatorname{sech}^2(v+\sqrt qZ)$ for $0\le v\le1$
and $q\le1$, restricting to $|Z|\le1$; larger means follow by
monotonicity. Choose $A=2/c_0$. Repeating this shift doubles the
minimum signed mean until it reaches $c_0/2$, and preserves this
lower bound thereafter.

Choose
\[
 m=\max\{0,\lceil\log_2(c_0/(2u_*))\rceil\},\qquad
 H=L-m-3.
\]
The stated depth condition, with a sufficiently large $C_1$, ensures
$H\ge0$. Use one initial seed layer, $H$ free layers, $m$ shift
layers, and one final saturation layer. This is exactly $L-1$ hidden
layers and avoids any circular choice of the cooling time.
For the final shift use $B b$, with $B$ large enough in terms of
$\eta$ that its feature mean $v$ obeys $|v_i-y_i|<\eta$.
The bound
\[
 1-y_i\mathbb E\tanh(G_i+B b_i)\le2e^{2-Bc_0}
\]
suffices. Its entropy cost is at most $B^2/2$. Finally $Q_{L-1}\succeq vv^T$,
so $z=av$ has readout cost at most $a^2/2$.
The total cost is at most $1+mA^2/2+B^2/2+a^2/2$, which has the
asserted bound. Every Gaussian shift and rank-one domination is on
the full finite domain throughout the construction.
\end{proof}

\subsection{Comparison of posterior masses}

Correct prediction at both members of a test pair has a larger prior
rate than fitting the training sample. Dividing the former probability
by the latter therefore makes simultaneous correctness on a fixed pair
vanish as width grows. Averaging over the pairing then gives the risk
and confidence bounds for the full posterior.

\begin{proof}[Proof of Theorem~\ref{thm:critical-regression-depth-lower}]
By Lemma~\ref{lem:separation}, the training scalar coordinates are
$\Delta_n$-separated with probability at least $49/50$, where
$\Delta_n=(100n^2)^{-1/\alpha}$ and $\alpha=\log2/\log20$.
Fix any such sample. Since the Cantor measure is invariant under $F$,
every teacher label belongs to
$[-13/20,-1/2]\cup[1/2,13/20]$.
Set $y_i=\operatorname{sign}f_h(X_i)$ and $a=3/5$. Thus
\[
 |f_h(X_i)-ay_i|\le1/10.
\]
Apply Lemma~\ref{lem:critical-mean-seed-depth} on the finite training
set. Its quantity $D/m_n^2$ is polynomial in $n$.
Lemma~\ref{lem:critical-late-amplification-depth}, with $\eta=1/100$,
then supplies an output whose training errors are less than
\[
 1/10+a\eta=53/500<1/8.
\]
Reduce the universal $c$ and increase $h_0$ so that, whenever
$n\le e^{ch}$, the depth condition holds and this witness's rate is
less than $h/200$. The open training event $A_S$ therefore satisfies
\begin{equation}\label{eq:critical-training-rate-depth}
 \liminf_{N\to\infty}\frac1N\log\Pr(A_S)>-h/200.
\end{equation}
The strict inequality follows from the strict witness cost bound.

Now fix one test pair $(X_+,X_-)$ in the balanced pairing of
Lemma~\ref{lem:pairing}. Let $B_{X_+,X_-}$ be the closed event that
both real outputs are within $1/4$ of their teacher values.
Apply the two-point version of Lemma~\ref{lem:critical-finite-ldp-depth}.
Equation~\eqref{eq:critical-correct-pair-rate} gives
\[
 \limsup_{N\to\infty}\frac1N\log\Pr(B_{X_+,X_-})\le-h/100.
\]
Together with~\eqref{eq:critical-training-rate-depth}, the elementary
ratio bound
\[
 \Pr_{w\mid S}(B_{X_+,X_-})
 \le\frac{\Pr(B_{X_+,X_-})}{\Pr(A_S)}
\]
shows that this posterior probability tends to zero for every fixed
test pair. No joint or functional large-deviation principle is being
assumed. The width at which the ratio becomes small may depend on
the test pair.

Let $b_N(X_+,X_-)$ denote this pair-correct probability and average
with respect to the pairing law. Bounded convergence gives
$\mathbb E b_N\to0$. Inclusion--exclusion gives, pair by pair,
\[
 p_N(X_+)+p_N(X_-)\le1+b_N(X_+,X_-).
\]
A fair choice between the paired inputs has law $D$. Averaging proves
that the posterior Gibbs error is at least
$(1-\mathbb E b_N)/2$, giving its asserted lower limit.

For the confidence assertion, if $b_N<1/3$, both correct-output
probabilities cannot be at least $2/3$. At least one member of the
pair has probability strictly below $2/3$. Consequently
\[
 D\{x:p_N(x)<2/3\}\ge\frac12\Pr_{\rm pair}\{b_N<1/3\}
 \ge\frac{1-3\mathbb E b_N}{2}.
\]
Taking lower limits proves the confidence claim, since $\mathbb E b_N\to0$.
The argument applies
for every separated sample and handles arbitrary posterior mixtures.
\end{proof}
